So far, the control methods that we have studied, in particular LQR and Model Predictive Control, have relied on the availability of a state-space model of the plant. In discrete time, this is typically written as
\[
x_{t+1} = A x_t + B u_t, 
\qquad
y_t = C x_t,
\]
or, more generally, with a nonlinear update map. This representation is extremely convenient for control design, because it provides a Markovian description of the system: once the current state is known, the future evolution depends only on the present state and on the input that will be applied. In this sense, the state is the variable that summarizes all the information from the past that is relevant for predicting the future.

Having such a model, however, means having a very detailed description of the plant. One must know the parameters that govern the state update, the way in which the input affects the state evolution, and the way in which the internal state is mapped into the measured outputs. Even more fundamentally, one must decide what the state actually is, that is, which variables are sufficient to obtain a Markovian model. This is already a nontrivial modeling choice, especially when the internal physical variables are not directly measurable.

A first important observation is that a state-space representation is not unique. Different internal coordinates may describe exactly the same external input-output behavior. Indeed, if we introduce a new state variable \(w_t\) through an invertible linear change of coordinates
\[
x_t = T w_t,
\]
with \(T\) nonsingular, then the system can be equivalently written as
\[
w_{t+1} = T^{-1} A T w_t + T^{-1} B u_t,
\qquad
y_t = C T w_t.
\]
Although the matrices appearing in the model have changed, the input-output trajectories remain exactly the same. Therefore, the state should not be interpreted as a uniquely defined physical object, but rather as an internal variable that makes prediction and control possible. What really matters from the external viewpoint is the input-output behavior of the system.

This point is one of the main motivations for system identification. If different state-space models can generate the same measured behavior, then reconstructing a model from data cannot be understood as recovering the ``true'' internal realization in an absolute sense. Instead, the goal is to build a model that reproduces the observed dynamics well enough for the intended task. In our context, this task is control. Therefore, the identified model should be sufficiently accurate to predict the future effect of the control inputs and to support the synthesis of model-based controllers.

System identification addresses exactly this problem. Rather than deriving the model entirely from first principles, we use measured input-output data to infer a dynamical description of the plant. This is especially useful when a first-principles model is unavailable, too expensive to derive, or too inaccurate for control purposes. The identified model may be a state-space realization, an input-output model, or another predictive representation, depending on the assumptions and on the control architecture of interest.

From this viewpoint, identification can be seen as the bridge between data and model-based control. On the one hand, advanced controllers such as LQR and MPC require a predictive model. On the other hand, the model need not be known a priori: it can be estimated from experiments. The central question of this chapter is therefore how to extract, from measured data, a model that captures the relevant dynamical behavior of the plant and can be reliably used for prediction and control.

A natural way to move from an internal state-space description to an external input-output description is to write the output as a function of the initial condition and of the past inputs. For the discrete-time linear system
\[
x_{t+1}=Ax_t+Bu_t,\qquad y_t=Cx_t,
\]
repeated substitution gives
\[
y_t = CA^t x_0 + \sum_{j=0}^{t-1} CA^{\,t-1-j} B\,u_j.
\]
This expression clearly separates two contributions. The term \(CA^t x_0\) is the response generated by the initial state, while the summation describes how past inputs are propagated to the output through the dynamics. The latter has the form of a discrete-time convolution and shows that, for linear time-invariant systems, the full input-output behavior is encoded in the sequence of matrices
\[
CA^{t-1}B,\qquad t>0.
\]

In the SISO case, where both \(u_t\in\mathbb{R}\) and \(y_t\in\mathbb{R}\), this sequence reduces to a scalar sequence
\[
h_t = CA^{t-1}B,\qquad t>0,
\]
which is the impulse response of the system. If the initial condition is zero, then the output produced by an arbitrary input sequence is completely determined by the convolution of the input with \(\{h_t\}_{t\geq 1}\). In this sense, the impulse response is a complete representation of the zero-state input-output map of the plant. It contains exactly the information needed to predict future outputs from past inputs, without explicitly referring to the internal state coordinates. This is especially important in identification, because the state itself is not directly unique, whereas the input-output behavior is the physically observable object.

The same idea extends immediately to MIMO systems. If the plant has \(m\) inputs and \(p\) outputs, then each coefficient
\[
H_t = CA^{t-1}B \in \mathbb{R}^{p\times m},\qquad t>0,
\]
is a matrix. Its entry in position \((i,j)\) represents the response of output \(i\) to a unit impulse applied on input \(j\). The sequence \(\{H_t\}_{t\geq 1}\) is called the sequence of Markov parameters of the system. These parameters generalize the impulse response and provide a compact description of the input-output behavior of the linear system. Once they are known, the output can be reconstructed from the past inputs through matrix convolution.

This viewpoint is also attractive from an experimental perspective. On a real plant, the Markov parameters can be estimated directly from data by exciting the system and measuring the corresponding outputs. In the ideal noise-free case, one can initialize the plant at zero and apply a unit impulse to one input channel at a time, while keeping the others equal to zero. For a MIMO system, this means that \(m\) experiments are sufficient in principle: each experiment isolates one column of the matrices \(H_t\), and by repeating the procedure for all inputs one reconstructs the full sequence of Markov parameters. In practice, however, perfect impulses cannot be generated and measurements are affected by noise, disturbances, and possible offsets in the initial condition. For this reason, identification methods usually rely on richer input signals and estimate the Markov parameters from batches of input-output data, rather than from a single idealized impulse experiment.

The key message is that, before identifying a state-space realization, one may already identify the external behavior of the plant through its impulse response, or equivalently through its Markov parameters. This provides the first bridge between measured data and predictive models, and it motivates the identification techniques developed in the rest of this chapter.

\subsection{Controllability, observability, and minimal realizations}

When the goal is to identify a state-space model from input-output data, not every realization is equally meaningful. Since the state is not unique, one can always introduce additional coordinates that do not change the external behavior of the system. For identification and control, however, we are interested in realizations in which the state has an actual dynamical role. This naturally leads to the notions of controllability and observability, which characterize whether the state can be influenced by the input and reconstructed from the output.

For the discrete-time linear system
\[
x_{t+1}=Ax_t+Bu_t,\qquad y_t=Cx_t,
\]
controllability concerns the possibility of steering the system from a given initial condition to a desired target state. To fix ideas, let us start from \(x_0=0\). After one step, the state is
\[
x_1 = Bu_0,
\]
so the set of states reachable in one step is exactly the image of \(B\). After two steps,
\[
x_2 = ABu_0 + Bu_1,
\]
hence the reachable set is the image of the matrix \([\,B\;\;AB\,]\). Continuing in the same way, after \(n\) steps the state belongs to the image of
\[
\mathcal{C} = \begin{bmatrix} B & AB & A^2B & \cdots & A^{n-1}B \end{bmatrix},
\]
which is called the controllability matrix. Therefore, the system is controllable precisely when every state in \(\mathbb{R}^n\) can be reached, that is, when
\[
\operatorname{rank}(\mathcal{C}) = n.
\]
This condition means that the columns generated by repeatedly propagating the input directions through the dynamics span the whole state space. It is worth stressing that \(n\) steps are enough to decide controllability: if the full state space has not been reached by then, taking more steps cannot generate new independent directions.

The controllability matrix also provides a constructive way to compute a control sequence that reaches a desired state. Indeed, for a target \(\bar x\) reached in \(n\) steps, one can write
\[
x_n = \mathcal{C}
\begin{bmatrix}
u_{n-1}\\
\vdots\\
u_0
\end{bmatrix}.
\]
If \(\mathcal{C}\) is square and invertible, which is possible in particular in the scalar-input case, the input sequence is obtained directly by inversion. In the more general case, \(\mathcal{C}\) is not square, and the problem may admit multiple solutions. A standard choice is then the minimum-norm solution, which can be computed through the pseudoinverse.

Dual to controllability is observability. While controllability asks whether the input can excite all internal directions, observability asks whether the effect of the initial state can be detected from the output. Starting again from the linear model, at time \(t=0\) one has
\[
y_0 = Cx_0,
\]
so any initial state in the kernel of \(C\) is invisible at the output at that instant. After one step,
\[
y_1 = CAx_0,
\]
hence a state that belongs simultaneously to the kernels of \(C\) and \(CA\) remains indistinguishable over the first two output samples. Proceeding recursively, after \(n\) steps the invisible states are those in the kernel of
\[
\mathcal{O}=
\begin{bmatrix}
C\\
CA\\
CA^2\\
\vdots\\
CA^{n-1}
\end{bmatrix},
\]
which is called the observability matrix. The system is observable if and only if
\[
\operatorname{rank}(\mathcal{O})=n.
\]
Equivalently, distinct initial states always generate distinct output trajectories over a finite observation window. Also here, \(n\) steps are sufficient: if a state cannot be distinguished from the output sequence of length \(n\), it will remain unobservable even if more measurements are collected.

As in the controllability case, observability has a constructive interpretation. Given an output sequence \(y_0,\dots,y_{n-1}\), one can write
\[
\begin{bmatrix}
y_0\\
y_1\\
\vdots\\
y_{n-1}
\end{bmatrix}
=
\mathcal{O}x_0.
\]
If \(\mathcal{O}\) is invertible, the initial state is uniquely determined. More generally, when \(\mathcal{O}\) has full column rank, the state can still be reconstructed uniquely, for instance through a least-squares inversion based on the pseudoinverse.

These two notions are central in system identification because they single out the realizations that are meaningful from input-output data. If a realization contains uncontrollable states, part of the state cannot be influenced by the input and is therefore irrelevant for prediction and control. If it contains unobservable states, part of the state has no effect on the measured output and therefore cannot be inferred from data. In both cases, the realization contains redundant internal coordinates. For this reason, in identification we focus on realizations that are both controllable and observable. Such realizations are called minimal, because they do not contain superfluous states and represent the input-output behavior with the smallest possible state dimension.

A simple example illustrates the point. Consider a plant whose output is a moving average of the past three inputs,
\[
y_t = a_1u_{t-1}+a_2u_{t-2}+a_3u_{t-3}.
\]
A natural state is obtained by storing delayed copies of the input:
\[
x_t=
\begin{bmatrix}
u_{t-1}\\
u_{t-2}\\
u_{t-3}
\end{bmatrix}.
\]
With this choice, the dynamics become
\[
x_{t+1}
=
\begin{bmatrix}
0&0&0\\
1&0&0\\
0&1&0
\end{bmatrix}x_t
+
\begin{bmatrix}
1\\
0\\
0
\end{bmatrix}u_t,
\qquad
y_t=
\begin{bmatrix}
a_1 & a_2 & a_3
\end{bmatrix}x_t.
\]
This realization has an immediate interpretation: the state is simply a memory of the past input values needed to compute the current output.

The controllability matrix of this realization is
\[
\mathcal{C}=
\begin{bmatrix}
B & AB & A^2B
\end{bmatrix}=I,
\]
so the system is controllable. This is consistent with the physical meaning of the state, since the input directly writes the first memory cell and, through the shift dynamics, can generate any delayed input profile. The observability matrix is
\[
\mathcal{O}=
\begin{bmatrix}
a_1 & a_2 & a_3\\
a_2 & a_3 & 0\\
a_3 & 0 & 0
\end{bmatrix}.
\]
Hence observability depends on the coefficients \(a_1,a_2,a_3\). When this matrix has full rank, the three stored delays all leave a visible signature in the output and the realization is minimal. If instead the rank is smaller, then some component of the chosen state does not affect the output independently, which means that a lower-order realization exists.

This example clarifies why controllability and observability are so important in identification. Starting from input-output data, one does not want to recover an arbitrary realization, but rather a realization whose state is both reachable from the input and visible from the output. Only in this case does the state provide a well-posed internal representation of the external behavior of the plant, and only in this case can we hope to obtain a compact and useful model for control design.

\subsection{Realization from data and the Ho--Kalman construction}

The discussion so far has shown that the input-output behavior of a linear time-invariant system is completely described by its Markov parameters
\[
\{CB,\;CAB,\;CA^2B,\;\dots\}.
\]
This naturally leads to the realization problem: if these parameters are known from experiments, how can one reconstruct a state-space model
\[
x_{t+1}=Ax_t+Bu_t,\qquad y_t=Cx_t
\]
that generates them? More generally, given generic input-output data, how can one construct a realization \(A,B,C\) consistent with the observed dynamics? The goal is not to recover a unique internal state, since different state-space models may have the same external behavior, but rather to compute a minimal realization, that is, one that is both controllable and observable.

The key object in this construction is the Hankel matrix built from the Markov parameters. If we define the observability and controllability matrices of a minimal realization as
\[
\mathcal{O}=
\begin{bmatrix}
C\\
CA\\
CA^2\\
\vdots\\
CA^{n-1}
\end{bmatrix},
\qquad
\mathcal{C}=
\begin{bmatrix}
B & AB & A^2B & \cdots & A^{n-1}B
\end{bmatrix},
\]
then their product is
\[
\mathcal{H}=\mathcal{O}\mathcal{C}.
\]
Expanding this product shows that every block of \(\mathcal{H}\) is one of the Markov parameters:
\[
\mathcal{H}=
\begin{bmatrix}
CB & CAB & CA^2B & \cdots & CA^{n-1}B\\
CAB & CA^2B & CA^3B & \cdots & CA^nB\\
CA^2B & CA^3B & CA^4B & \cdots & CA^{n+1}B\\
\vdots & \vdots & \vdots & \ddots & \vdots\\
CA^{n-1}B & CA^nB & CA^{n+1}B & \cdots & CA^{2n-2}B
\end{bmatrix}.
\]
Thus the Hankel matrix can be assembled directly from experimental data as soon as enough Markov parameters are available. This is the crucial bridge between input-output measurements and state-space realizations.

To make the structure visually apparent, it is often helpful to highlight repeated Markov parameters along the anti-diagonals. For example, one may write
\[
\mathcal{H}=
\left[
\begin{array}{ccccc}
\colorbox{red!20}{$CB$} &
\colorbox{blue!20}{$CAB$} &
\colorbox{green!20}{$CA^2B$} &
\cdots &
CA^{n-1}B\\[0.3em]
\colorbox{blue!20}{$CAB$} &
\colorbox{green!20}{$CA^2B$} &
CA^3B &
\cdots &
CA^nB\\[0.3em]
\colorbox{green!20}{$CA^2B$} &
CA^3B &
CA^4B &
\cdots &
CA^{n+1}B\\
\vdots & \vdots & \vdots & \ddots & \vdots\\
CA^{n-1}B & CA^nB & CA^{n+1}B & \cdots & CA^{2n-2}B
\end{array}
\right].
\]
The repeated colors emphasize that each anti-diagonal contains the same Markov parameter. This is exactly the signature of a Hankel structure: the matrix is constant along anti-diagonals, and those constants are the impulse-response coefficients of the system.

The Ho--Kalman algorithm uses this structure in three conceptual steps. First, one constructs the Hankel matrix from the measured Markov parameters. Second, one factorizes the Hankel matrix as
\[
\mathcal{H}=\mathcal{O}\mathcal{C}.
\]
Third, one extracts a realization \(A,B,C\) from the factors \(\mathcal{O}\) and \(\mathcal{C}\).

At first sight, the first step raises a basic difficulty: how many Markov parameters are needed? The answer depends on the order \(n\) of the minimal realization, which is unknown a priori. However, if the realization is minimal, then both \(\mathcal{O}\) and \(\mathcal{C}\) have rank \(n\), and therefore the Hankel matrix also has rank \(n\). This means that, in principle, one can enlarge the Hankel matrix progressively until its rank stops increasing. That plateau reveals the order of the system. In noise-free settings this idea works exactly; in practice, with noisy data, the same principle survives in approximate form and the effective rank is inferred numerically.

This rank-based interpretation is easy to understand on simple examples. A first-order low-pass filter has state-space model
\[
x_{t+1}=ax_t+u_t,\qquad y_t=x_t,
\]
so its Markov parameters are \(1,a,a^2,\dots\), and the corresponding Hankel matrix has rank one. By contrast, a frictionless cart model written as
\[
v_{t+1}=v_t+u_t,\qquad z_{t+1}=z_t+v_t,\qquad y_t=z_t
\]
has a two-dimensional state, and the Hankel matrix reflects this higher order through rank two. In this way, the rank of the Hankel matrix captures the intrinsic dimension of the minimal realization.

The second step of the Ho--Kalman method is the factorization
\[
\mathcal{H}=\mathcal{O}\mathcal{C}.
\]
An important conceptual point is that this factorization is not unique, but this is not a problem. Indeed, nonuniqueness of the factorization simply reflects nonuniqueness of the state coordinates. If we apply a change of basis
\[
\tilde A = T^{-1}AT,\qquad \tilde B = T^{-1}B,\qquad \tilde C = CT,
\]
with \(T\) invertible, then the corresponding observability and controllability matrices become
\[
\tilde{\mathcal{O}}=
\begin{bmatrix}
\tilde C\\
\tilde C\tilde A\\
\tilde C\tilde A^2\\
\vdots
\end{bmatrix}
=
\begin{bmatrix}
CT\\
CTA T^{-1}T\\
CTA^2T^{-1}T\\
\vdots
\end{bmatrix}
=
\mathcal{O}T,
\]
and similarly
\[
\tilde{\mathcal{C}}=T^{-1}\mathcal{C}.
\]
Therefore,
\[
\tilde{\mathcal{O}}\tilde{\mathcal{C}}
=
\mathcal{O}TT^{-1}\mathcal{C}
=
\mathcal{O}\mathcal{C}
=
\mathcal{H}.
\]
So different factorizations of \(\mathcal{H}\) correspond exactly to different coordinate representations of the same input-output system. The factorization does not identify a unique realization, but it identifies the system up to similarity transformation, which is the correct notion of uniqueness for state-space models.

Because of this, the factorization step leaves some freedom. In principle, one could take trivial choices such as \(\mathcal{O}=I\) and \(\mathcal{C}=\mathcal{H}\), or \(\mathcal{O}=\mathcal{H}\) and \(\mathcal{C}=I\), but these choices are not useful numerically and do not directly reveal a well-scaled state representation. In practice, one prefers factorizations based on the singular value decomposition, because they expose the rank of the Hankel matrix and provide numerically robust coordinates for the state.

Once a factorization \(\mathcal{H}=\mathcal{O}\mathcal{C}\) has been obtained, the realization matrices are recovered from the shift structure of \(\mathcal{O}\) and \(\mathcal{C}\). The matrix \(C\) is simply the first block row of \(\mathcal{O}\), while \(B\) is the first block column of \(\mathcal{C}\). The state transition matrix \(A\) is then determined by exploiting the fact that consecutive block rows of \(\mathcal{O}\) differ by multiplication by \(A\), or equivalently that consecutive block columns of \(\mathcal{C}\) differ by the same shift. In this way, the whole state-space realization is reconstructed from the Hankel matrix, and therefore from measured input-output behavior alone.

The significance of the Ho--Kalman construction is that it turns impulse-response data into a minimal state-space model through linear-algebraic operations. It provides one of the cleanest answers to the realization problem: first identify the Markov parameters, then organize them into a Hankel matrix, then factorize this matrix to uncover the underlying controllable and observable state structure.

\subsection{SVD-based Ho--Kalman realization}

The key missing ingredient in the Ho--Kalman construction is a numerically reliable way to factorize the Hankel matrix. This is provided by the singular value decomposition. If \(H\in\mathbb{R}^{m\times n}\) has rank \(r\), then it can be written as
\[
H = U\Sigma V^\top,
\]
where \(U\in\mathbb{R}^{m\times m}\) and \(V\in\mathbb{R}^{n\times n}\) are orthogonal matrices, while \(\Sigma\in\mathbb{R}^{m\times n}\) is diagonal in the rectangular sense, with nonnegative singular values on the diagonal. If the rank of \(H\) is \(r\), then only the first \(r\) singular values are nonzero, and the decomposition can be partitioned as
\[
H=
\begin{bmatrix}
U_1 & U_2
\end{bmatrix}
\begin{bmatrix}
\Sigma_1 & 0\\
0 & 0
\end{bmatrix}
\begin{bmatrix}
V_1 & V_2
\end{bmatrix}^{\!\top}
=
U_1\Sigma_1V_1^\top,
\]
where \(\Sigma_1\in\mathbb{R}^{r\times r}\) contains the nonzero singular values. This decomposition is one of the most useful tools in numerical linear algebra: it is computationally efficient, numerically robust, and it reveals directly the rank and the dominant subspaces of the matrix. In particular, the columns of \(U_1\) span the image of \(H\), while the dimension of \(\Sigma_1\) identifies the rank.

This is exactly the type of factorization we need for the Hankel matrix. Instead of building a square Hankel matrix of unknown dimension, it is convenient to introduce an extended Hankel matrix
\[
\mathcal{H}_{k,\ell}=
\begin{bmatrix}
G_1 & G_2 & \cdots & G_\ell\\
G_2 & G_3 & \cdots & G_{\ell+1}\\
G_3 & G_4 & \cdots & G_{\ell+2}\\
\vdots & \vdots & \ddots & \vdots\\
G_k & G_{k+1} & \cdots & G_{k+\ell-1}
\end{bmatrix},
\]
where \(G_t = CA^{t-1}B\) denotes the \(t\)-th Markov parameter. Here \(k\) and \(\ell\) are chosen larger than the unknown system order \(n\). If the data are generated by an \(n\)-dimensional minimal realization, then this extended Hankel matrix factors as
\[
\mathcal{H}_{k,\ell}=\mathcal{O}_k\mathcal{C}_\ell,
\]
with
\[
\mathcal{O}_k=
\begin{bmatrix}
C\\
CA\\
CA^2\\
\vdots\\
CA^{k-1}
\end{bmatrix},
\qquad
\mathcal{C}_\ell=
\begin{bmatrix}
B & AB & A^2B & \cdots & A^{\ell-1}B
\end{bmatrix}.
\]
These are the extended observability and controllability matrices. As soon as \(k\ge n\) and \(\ell\ge n\), both matrices have rank \(n\), and therefore \(\mathcal{H}_{k,\ell}\) also has rank \(n\). This is the fundamental reason why one chooses \(k\) and \(\ell\) larger than the system order: once the horizon is long enough, the rank of the Hankel matrix reveals the true dimension of the minimal realization.

At this point, the SVD gives a natural and balanced factorization of \(\mathcal{H}_{k,\ell}\). If
\[
\mathcal{H}_{k,\ell}=U_1\Sigma_1V_1^\top,
\]
where \(\Sigma_1\) contains the nonzero singular values, then the order of the system is simply
\[
n=\operatorname{rank}(\mathcal{H}_{k,\ell})=\dim(\Sigma_1).
\]
Moreover, we may decompose the Hankel matrix as
\[
\mathcal{H}_{k,\ell}
=
\underbrace{U_1\Sigma_1^{1/2}}_{\mathcal{O}_k}
\underbrace{\Sigma_1^{1/2}V_1^\top}_{\mathcal{C}_\ell}.
\]
This choice is especially convenient because it distributes the scaling symmetrically between the observability and controllability factors. In this way we obtain numerical representatives of the extended observability and controllability matrices:
\[
\mathcal{O}_k = U_1\Sigma_1^{1/2},
\qquad
\mathcal{C}_\ell = \Sigma_1^{1/2}V_1^\top.
\]

From these matrices, \(C\) and \(B\) are obtained immediately. Indeed, \(C\) is the first block row of \(\mathcal{O}_k\), while \(B\) is the first block column of \(\mathcal{C}_\ell\). Visually, one can emphasize this as
\[
\mathcal{O}_k=
\begin{bmatrix}
\colorbox{green!20}{$C$}\\
CA\\
CA^2\\
\vdots\\
CA^{k-1}
\end{bmatrix},
\qquad
\mathcal{C}_\ell=
\begin{bmatrix}
\colorbox{green!20}{$B$} & AB & A^2B & \cdots & A^{\ell-1}B
\end{bmatrix}.
\]
Thus the first visible block of the observability factor gives the output matrix, and the first visible block of the controllability factor gives the input matrix.

To recover the state transition matrix \(A\), one introduces the shifted Hankel matrix
\[
\mathcal{H}^{\uparrow}_{k,\ell}=
\begin{bmatrix}
G_2 & G_3 & \cdots & G_{\ell+1}\\
G_3 & G_4 & \cdots & G_{\ell+2}\\
G_4 & G_5 & \cdots & G_{\ell+3}\\
\vdots & \vdots & \ddots & \vdots\\
G_{k+1} & G_{k+2} & \cdots & G_{k+\ell}
\end{bmatrix}.
\]
By construction, this matrix satisfies
\[
\mathcal{H}^{\uparrow}_{k,\ell}=
\begin{bmatrix}
CAB & CA^2B & CA^3B & \cdots\\
CA^2B & CA^3B & CA^4B & \cdots\\
CA^3B & CA^4B & CA^5B & \cdots\\
\vdots & \vdots & \vdots & \ddots
\end{bmatrix}
=
\mathcal{O}_k A \mathcal{C}_\ell.
\]
Once \(\mathcal{O}_k\) and \(\mathcal{C}_\ell\) are known, this relation can be inverted in the least-squares sense, yielding
\[
A = \mathcal{O}_k^\dagger \mathcal{H}^{\uparrow}_{k,\ell}\mathcal{C}_\ell^\dagger,
\]
where \({}^\dagger\) denotes the pseudoinverse. Therefore the full realization \((A,B,C)\) is reconstructed directly from the measured Markov parameters.

Putting everything together, the Ho--Kalman procedure can be summarized as follows. First, one collects impulse-response data and estimates the Markov parameters of the system. Then one chooses \(k,\ell\) sufficiently large and builds the extended Hankel matrix \(\mathcal{H}_{k,\ell}\) together with its shifted version \(\mathcal{H}^{\uparrow}_{k,\ell}\). Next, one computes the SVD of \(\mathcal{H}_{k,\ell}\), uses the nonzero singular values to determine the system order \(n\), and defines the extended observability and controllability matrices through the balanced factorization
\[
\mathcal{O}_k = U_1\Sigma_1^{1/2},
\qquad
\mathcal{C}_\ell = \Sigma_1^{1/2}V_1^\top.
\]
Finally, one extracts
\[
C=\text{first block row of }\mathcal{O}_k,\qquad
B=\text{first block column of }\mathcal{C}_\ell,
\]
and computes
\[
A = \mathcal{O}_k^\dagger \mathcal{H}^{\uparrow}_{k,\ell}\mathcal{C}_\ell^\dagger.
\]

The main message is that, even when a state-space model is not available a priori, it can be reconstructed from data. In this sense, the data themselves contain the model information needed for prediction and control. Once a realization has been identified, the model-based control methods developed in the previous chapters, such as LQR and MPC, can be applied to the identified plant exactly as if the model had been known from first principles.

At the same time, it is important to keep in mind the assumptions underlying this construction. The derivation assumes that the plant is an exact finite-dimensional linear time-invariant system and that the Markov parameters are known without noise. Under these ideal conditions, the rank of the Hankel matrix is sharp, the realization is exact, and the Ho--Kalman algorithm recovers a minimal state-space model up to a change of coordinates. Real identification problems are usually more challenging: measurements are noisy, data are finite, and the true dynamics may be nonlinear or only approximately finite-dimensional. Therefore, the material presented here should be understood as a first and fundamental example of system identification, rather than as a complete treatment of the subject. More advanced settings include noisy and stochastic data, missing measurements, prior structural information, nonlinear models, reduced-order representations, and both parametric and nonparametric identification methods. Still, the ideas developed in this chapter remain central: organize data into structured matrices, reveal the intrinsic dimension through linear algebra, and reconstruct a model that captures the dynamical behavior relevant for control.